% !TeX spellcheck = it_IT
\newpage
\section{Introduzione}
L'obiettivo del corso è di studiare gli strumenti necessari a comprendere e gestire il processo di sviluppo delle interfacce e dei prodotti interattivi. È fondamentale che l'informatico crei un prodotto \textbf{apprezzato} dalle persone.

\subsection{Progettazione}
\subsubsection{Design}
\begin{definition}[Design]
	Per design si intende sia il processo di \textbf{progettazione} e \textbf{pianificazione} sia l'\textbf{output} stesso di questo processo.
\end{definition}
Nel design il primo passo è capire \textbf{perché} un problema esiste e se ha davvero bisogno di una soluzione. È quindi fondamentale fare in modo che il software risolva i problemi dell'utente, senza aggiungere funzionalità inutili non necessarie all'utente.\\

Nel design del prodotto è spesso necessario modificare requirement e specifiche per andare in contro alle esigenze degli utenti, rendendo quindi lo sviluppo un processo \textbf{antropocentrico} di adattamento del software.\\\\
Esistono varie sotto discipline del design, tra cui:
\begin{itemize}
	\item \textbf{Interaction} design: l'attività di progettazione dell'interazione che avviene tra esseri umani e oggetti in generale. Ha come obiettivo quello di rendere macchine, servizi e sistemi \textbf{usabili} per gli utenti target, mettendo al centro i loro \textbf{bisogni}. Si suddivide in:
	\begin{itemize}
		\item Design di \textbf{prodotto}: progettazione di beni e servizi il cui obiettivo principale è quello di essere utilizzati da quanti più utenti possibili migliorandone la vita\footnote{Spesso usato in modo interscambiabile con \textbf{design industriale}: processo strategico di risoluzione dei problemi che guida l'innovazione e porta a una migliore qualità della vita attraverso prodotti, sistemi, servizi ed esperienze innovative}
		\begin{definition}[Designer]
			L'inventore o designer del prodotto si focalizza su un problema da risolvere e inventa un prodotto che grazie alle sue caratteristiche tecniche permette di risolvere un problema.
		\end{definition}
		\item Design dell'\textbf{esperienza utente} o \textbf{UX}: il processo volto ad aumentare la soddisfazione e la fedeltà del cliente migliorando l'\textbf{usabilità} e la piacevolezza. Aggiunge alla tradizionale progettazione gli aspetti di business, marketing e sviluppo prodotto.
		\begin{definition}[UX Designer]
			Lo UX Designer ha l'obiettivo di far vivere all'utente del suo prodotto la miglior esperienza possibile, evitando di farlo sentire frustrato e deluso.
		\end{definition}
		\item Design dell'\textbf{interfaccia} o \textbf{UI}: ha come focus il modo in cui le persone interagiscono con la tecnologia e cerca di migliorare la loro comprensione di ciò che si può fare. Dalla UX si crea uno schema di interazione che serve a progettare l'interfaccia utente al fine di abilitare l'esperienza progettata.
		\begin{definition}[UI Designer]
			Lo UI Designer progetta l'aspetto \textbf{estetico} e la struttura dell'interfaccia così che questa durante l'utilizzo induca l'utente nel seguire l'esperienza progettata. Produce un \textbf{wireframe}\footnote{Bozza grafica} dell'interfaccia e delle linee guida che verranno seguite dagli \textbf{UI Developer} per l'implementazione.
		\end{definition}
	\end{itemize}
	\item \textbf{Human Machine} e \textbf{Human Computer} Interaction: ha come obiettivo quello di rendere possibile e \textbf{facilitare} al massimo per un essere umano l'\textbf{uso} e l'\textbf{interazione} con i sistemi del mondo IT
\end{itemize}

\begin{center}
	\includegraphics[scale=.5]{UXUI_FlowDiagram_800.jpg}
\end{center}

\subsubsection{Pensiero computazionale}
Nel mondo del \textit{design}, in particolare quello industriale, ci si approccia alla progettazione e risoluzione dei problemi in modo diverso dal mondo informatico.\\
Per questo motivo nel 2006 Jeannette Wing, direttrice del dipartimento di Informatica della Mellon University, formula la seguente definizione:
\begin{definition}[Pensiero computazionale]
	Il pensiero computazionale è un processo di formulazione di problemi e soluzioni in una forma che sia eseguibile da un agente che processi informazioni.
\end{definition}
Questo consiste non solo nel saper programmare ma anche nel pensare a diversi \textbf{livelli di astrazione}, abituando al \textit{rigore} e rendendo possibili gli \textit{atti creativi}. Permette di interagire con persone e strumenti e di usare le macchine come oggetti capaci di compensare le lentezze o l'imprecisione dell'uomo.
\subsubsection{Confronto}
Se nel \textbf{design} la progettazione è affrontata in maniera globale, nel tentativo di comprendere il problema nel suo insieme, nel \textbf{pensiero computazionale} si basa sulla suddivisione di un problema in sotto-problemi in modo da giungere ad una formulazione del problema come \textit{algoritmo}.\\

È importante evidenziare che il design dell'interazione e il pensiero computazionale non sono mutualmente esclusivi ma anzi è nell'\textbf{unione} dei due e nella loro \textbf{integrazione} che si crea software di qualità.